\subsection{Indirect Support Prioritizes Drug Targets With Higher Relative Success}

\begin{figure}[t!]
    \centering
    {{canva:5}}
    \caption{Indirect and combined PIGEAN genetic support enrich for successful drug targets compared with Open Targets metrics.}
    \label{fig:validation-pharma}
\end{figure}

% Paragraph on continuous drug target relative success and comparing to Open Targets
We evaluated drug targets prioritized by PIGEAN's indirect genetic support, defining a target as a gene and semantically matched indication pair (Methods) with indirect support above a specified log-odds threshold (likelihood of relevance) across {{NUM_TOTAL_TRAITS_MOUSE_MSIGDB}} phenotypes, comprising {{NUM_COMMON_TRAITS_MOUSE_MSIGDB}} common traits and {{NUM_RARE_TRAITS_MOUSE_MSIGDB}} rare diseases. At the relative-success (RS) level achieved by Open Targets Genetics (OTG) with a minimum locus-to-gene (L2G) score $\geq 0.1$ as curated by Minikel et al.\ \cite{minikel_refining_2024} (RS = {{MINIKEL_OTG_RS:.2f}}; 95\% CI: {{MINIKEL_OTG_RS_CI_LOWER:.2f}}--{{MINIKEL_OTG_RS_CI_UPPER:.2f}}), PIGEAN's indirect support prioritized {{NUM_PIGEAN_PRIORITIZED_TARGETS_MOUSE_MSIGDB_AT_RS_eq_MINIKEL_OTG_INDIRECT}} targets, compared with only {{NUM_MINIKEL_OTG_PRIORITIZED_TARGETS}} targets prioritized by OTG. When direct and indirect genetic support were combined (combined genetic support) at the same RS level {{MINIKEL_OTG_RS}}, PIGEAN prioritized {{NUM_PIGEAN_PRIORITIZED_TARGETS_MOUSE_MSIGDB_AT_RS_eq_MINIKEL_OTG_COMBINED}} targets, but excluded {{NUM_UNIQUE_PIGEAN_INDIRECT_TARGETS_COMPARED_WITH_COMBINED}} targets uniquely identified by indirect support (with {{NUM_SHARED_PIGEAN_TARGETS_INDIRECT_AND_DIRECT}} targets shared between the indirect-only and combined sets), indicating that indirect and combined support surface partially distinct sets of drug targets.

Because OTG and its corresponding drug discovery resource, the Open Targets Platform (OTP), are the only settings in which a continuous RS metric is available, we benchmarked PIGEAN's indirect and combined genetic support against the Minikel et al.\ curation of OTG targets as a function of L2G score, as well as against OTP's harmonic mean score for genetic support and its literature- and model-based concept of indirect genetic support (Methods; Figure~\ref{fig:validation-pharma}a). RS was plotted as a function of the fraction of each resource prioritized at a given threshold of its associated score, parameterized by a threshold $\alpha \in [0,1]$. Across this range, PIGEAN's indirect and combined support consistently yielded higher RS than OTP's indirect-support and all-genetic-support metrics, while prioritizing a similar number of targets at $\alpha = 0.1$ (10\% of each resource). Although the Minikel et al.\ OTG curation performed slightly better than PIGEAN's indirect and combined support, this curated set was considerably smaller, containing only {{NUM_MINIKEL_OTG_TARGETS_AT_ALPHA_10}} targets at $\alpha = 0.1$.

We next examined OTG L2G score as a joint function of RS and the number of targets prioritized (Extended Data Figure~\ref{fig:extended-1}). L2G alone performed almost equivalently to PIGEAN's indirect genetic support, whereas combined PIGEAN support substantially outperformed all other approaches over the first 1,000 prioritized targets. Notably, RS for L2G remained relatively stable across the range of L2G scores, such that even scores close to 1 (nominally near-causal) did not yield markedly increased RS, consistent with the findings of Minikel et al. By contrast, both indirect and combined genetic support from PIGEAN achieved substantially higher RS at higher scores, indicating greater selectivity at more stringent thresholds.

% Paragraph on drug target success with temporal holdout of drug targets
Because the gene sets used by our PIGEAN models may contain latent information about drug targets, we further evaluated indirect and combined genetic support in a temporal holdout of drug targets over the previous 67 common traits with a GWAS published before 2018, using the 2018 versions of the MSigDB and mouse phenotype gene sets (same model as reported in Section 2.2). In this temporal analysis, we restricted evaluation to drug targets with launch dates after 2018. The Pharmaprojects dataset reports only the year of drug launch \cite{minikel_refining_2024} and does not capture the year a development program began, so we could exclude targets launched before 2018 but could not censor pre-2018 programs that did not launch; such programs were therefore retained in our RS calculations. Even under this temporal holdout, indirect and combined PIGEAN support remained strong predictors of RS for post-2018 drug targets (Figure~\ref{fig:validation-pharma}b), with indirect support achieving a maximum RS of {{MAX_RS_INDIRECT_TEMPORAL_HOLDOUT}} (95\% CI: {{MAX_RS_INDIRECT_TEMPORAL_HOLDOUT_CI_LOWER}}--{{MAX_RS_INDIRECT_TEMPORAL_HOLDOUT_CI_UPPER}}; $N$ = {{NUM_DRUG_TARGETS_SUPPORTED_INDIRECT_TEMPORAL_HOLDOUT}}) and combined support achieving a maximum RS of {{MAX_RS_COMBINED_TEMPORAL_HOLDOUT}} (95\% CI: {{MAX_RS_COMBINED_TEMPORAL_HOLDOUT_CI_LOWER}}--{{MAX_RS_COMBINED_TEMPORAL_HOLDOUT_CI_UPPER}}; $N$ = {{NUM_DRUG_TARGETS_SUPPORTED_COMBINED_TEMPORAL_HOLDOUT}}), with {{NUM_DRUG_TARGETS_SHARED_INDIRECT_AND_COMBINED_TEMPORAL_HOLDOUT}} targets shared between the indirect and combined sets. At $\alpha = 0.02$ (2\% of the resource), RS for indirect support stabilized at approximately {{RS_INDIRECT_TEMPORAL_HOLDOUT_STABLE_AT_ALPHA_002}} and RS for combined support at approximately {{RS_COMBINED_TEMPORAL_HOLDOUT_STABLE_AT_ALPHA_002}}, based on {{NUM_TARGETS_INDIRECT_TEMPORAL_HOLDOUT_AT_ALPHA_002}} and {{NUM_TARGETS_COMBINED_TEMPORAL_HOLDOUT_AT_ALPHA_002}} targets at this threshold, respectively.

% Paragraph on drug target success in comparision to minikel et al.